\subsection{Programación de audiencias e historia de sus declaraciones}\label{sec:diseno-audiencias}

La audiencia se modela como una raíz estable perteneciente a un expediente, con un tipo inmutable y revisiones sucesivas de su programación. Cada revisión conserva fecha con desfase, modalidad, sede o conexión declarada, nota opcional y referencias exactas a las fichas seleccionadas del directorio. Los estados \texttt{scheduled} y \texttt{cancelled} son organizativos: una fecha pasada no demuestra que la audiencia se celebró, y cancelar una cita no declara la nulidad de un acto procesal. No se impone una sola audiencia de cada tipo por expediente.

El catálogo operativo distingue inicial, intermedia, debate de juicio e individualización de sanciones y reparación del daño. Las dos primeras requieren, respectivamente, Investigación e Intermedia registradas; las dos últimas, Juicio. La individualización requiere además un antecedente de condena declarado y una versión documental exacta como soporte. La etapa no permite inferir ese antecedente. La compatibilidad técnica de tipo y etapa ordena la captura, sin certificar la procedencia jurídica de la audiencia.

Programar o reemplazar exige un perfil penal completo, estado administrativo activo y expectativas sobre las revisiones administrativa y procesal. La reprogramación agrega una revisión con motivo; no sobrescribe la anterior. Cancelar conserva los valores, participantes, soporte y contexto de programación de la base exacta, y registra por separado la administración vigente al confirmar. Por ello puede cancelarse una cita después de avanzar la etapa o de archivar una ficha vinculada, mientras el expediente permanezca abierto y el actor conserve permiso. El cierre administrativo impide también esa mutación.

La selección de participantes admite de cero a 32 fichas del mismo expediente. Una selección nueva exige una revisión vigente y activa; conservar una referencia ya vinculada mantiene el nombre, perfil e identidad de aquella revisión, incluso si el directorio cambió después. La cuenta autora sigue siendo distinta de las personas u órganos representados. Una lista vacía expresa selección no registrada y ninguna selección demuestra asistencia efectiva.

La agenda agrupa cabeceras actuales de los expedientes autorizados, con intervalo explícito, orden por instante UTC e identificador y paginación acotada. El filtrado se realiza en persistencia antes de paginar. Su proyección excluye notas, conexiones e identidad de participantes; estos datos requieren consultar el detalle autorizado. La agenda puede cambiar entre páginas y no representa una instantánea global ni un calendario judicial de días hábiles.

Esta base cubre la organización de citas dentro del alcance de CU-08 y RF-07, pero su aceptación conserva los resultados, asistentes, acuerdos, plazos derivados y alertas pendientes. Los cuatro tipos ordinarios no agotan las menciones del marco teórico: la audiencia de medidas cautelares requiere tratamiento propio, y la continuación de audiencia requiere precisar su relación con la sesión original. El calendario unificado con plazos y sus vistas tampoco se deduce de una agenda de audiencias. La distinción mantiene parciales RF-07, RF-10 y el objetivo específico de gestión procesal.

{\raggedright El contrato se especifica en \rutarepo{docs/hearings-api.md}; la decisión de separar programación, historia y recibos se fundamenta en \rutarepo{docs/adr/0028-audited-hearing-scheduling.md}. La \cref{sec:impl-audiencias} desarrolla las fronteras de dominio, aplicación y transporte.\par}
